\documentclass[11pt,a4paper]{article}
\usepackage[margin=2.4cm]{geometry}
\usepackage{booktabs}
\usepackage{amsmath}
\usepackage{parskip}

\title{Is the price process non-stationary on monthly timescales?\\
\large The rolling-MLE drift study, and the verdict on periodic re-optimisation}
\author{Bayesian Capital}
\date{9 August 2026}

\begin{document}
\maketitle

\section*{The question}

The hypothesis under test: stock behaviour shifts with the AI-trade market view on
roughly monthly timescales, so parameters fitted to the full 2.2-year window are a stale
average, and a scheme that refits every one to three months on a trailing 3--12 month
window would track the market and out-earn the frozen configuration.

The refit scheme itself cannot be tested cheaply on P\&L --- its objective sees only
6--25 trades per short window against ten parameters --- so the hypothesis was tested
where it is identified: on prices. The Kalman local-linear-trend filter's noise scales
$(\lambda, \varphi_L, \psi)$ have an exact Gaussian likelihood through their one-step
innovations, giving 63--126 daily observations per window against three parameters.
If the market's regime changes the price process, these are the parameters of that
process, and the likelihood will see it move.

\section*{Method}

Rolling 63- and 126-day windows stepped 21 days across 2024-04-01 to 2026-08-07
(591 sessions), five names (RKLB, TSM, VST, VRT, MU). Per window: maximum-likelihood
fit of $(\lambda, \varphi_L, \psi)$ with standard errors from the Hessian, the OU
sleeve's rolling AR(1) coefficient and residual sigma at the deployed lookback, and
the AI factor's drawdown state (equal-weight TSM/VRT/VST/MU, as in the concentration
study). The instrument was validated on simulated data first: truth recovered within
2.5 standard errors in 7 of 8 126-day windows; median $|z|\approx 1$ on a stationary
path; a mid-sample doubling of $\varphi_L$ detected in 4 of 4 paths.

Two tests, deliberately independent of each other:
\begin{enumerate}
  \item \textbf{Error-bar test.} $\chi^2/\mathrm{dof}$ of non-overlapping window
        estimates about the full-sample MLE. Stationary calibration $\approx 1$.
  \item \textbf{Likelihood-ratio homogeneity test} (free of the Hessian): one vector
        for all windows vs one per window, $LR = 2(\sum_w \ell_w(\hat\theta_w) -
        \sum_w \ell_w(\hat\theta))$, $\chi^2_{3(K-1)}$. Plus the sharpest version of
        the hypothesis: the only AI drawdown episode in sample (January--April 2025,
        factor $-45.9\%$) as its own segment against the rest, 3 degrees of freedom.
\end{enumerate}

\section*{Results}

\textbf{The parameters do not move.} Every windows LR test accepts homogeneity, and
not narrowly:

\begin{center}
\begin{tabular}{lrrrrr}
\toprule
 & \multicolumn{2}{c}{63d windows ($K=9$)} & \multicolumn{2}{c}{126d windows ($K=4$)} & crash vs rest \\
 & LR (dof 24) & $p$ & LR (dof 9) & $p$ & $p$ (dof 3) \\
\midrule
MU   & 15.4 & 0.91 & 4.5 & 0.88 & 0.83 \\
RKLB & 15.6 & 0.90 & 4.3 & 0.89 & 0.19 \\
TSM  & 16.2 & 0.88 & 9.2 & 0.42 & 0.32 \\
VRT  & 18.7 & 0.77 & 7.5 & 0.59 & 0.04 \\
VST  & 13.3 & 0.96 & 2.7 & 0.98 & 0.81 \\
\bottomrule
\end{tabular}
\end{center}

The LR statistics sit \emph{below} their degrees of freedom throughout --- windows agree
with each other more closely than chance requires. Even the crash test fails: one name
in five reaches $p=0.04$, which across fifteen tests is what chance produces. The
Jan--Apr 2025 episode, the strongest regime change in the sample, did not change the
process parameters. The $\chi^2$/dof table agrees ($0.3$--$1.2$ almost everywhere;
TSM's $\lambda$ at 2.5 is the single mild exception), and $\psi$ is flat at
effectively zero in every window. The OU AR(1) coefficient is pinned at $0.95$--$0.98$
with IQR $0.02$--$0.06$ --- stable, and incidentally so close to a unit root that the
``mean-reversion'' forecast is a random walk anchor in practice, consistent with the
weekly finding that the MR term never binds.

\textbf{Why the intuition is right and the conclusion is wrong.} The market
\emph{does} change on monthly timescales --- the OU residual sigma moves from 2.0\% to
5.5\% of price across the sample. But the filter's noises are specified as
\emph{scales on the daily range} $F_t$: $q_L = (\varphi_L F_t)^2$, $r = (\lambda
F_t)^2$. The regime enters through $F_t$, day by day, with no refit --- faster than any
monthly re-optimisation could. What $(\lambda, \varphi_L, \psi)$ measure is the ratio
of level-shock, observation-noise and slope-noise to the range, and those ratios are
regime-invariant to within measurement error. The specification already contains the
adaptation the refit scheme was meant to add. The small coherent wobble that does
exist ($\varphi_L$ rising a few percent of its value in deep drawdowns, Spearman
$-0.4$ to $-0.7$ on overlapping windows) is second-order and inside the error bars.

\textbf{A side finding on what the P\&L optimiser fitted.} The full-sample MLE is far
from the deployed vector on every name --- e.g.\ TSM MLE
$\lambda/\varphi_L/\psi = 0.19/0.77/0.003$ against deployed $0.46/0.25/0.065$. The
likelihood says the close \emph{is} the level (small $\lambda$) and the level moves
$\sim\!0.6$--$0.8$ of the daily range; the deployed values instead buy heavy smoothing
--- a fair value that lags the market, producing deeper bids through $\mathrm{fair} -
k\sigma$. That is a \emph{policy} effect, not an estimate of the process, and it is
why the P\&L surface is a ridge: $k$ can compensate almost any $(\lambda, \varphi_L,
\psi)$. The deployed vector is not thereby wrong --- but these three parameters should
be understood as bid-shaping knobs, not statistics, and no amount of refitting them
will track a market regime.

\section*{Verdict}

Periodic re-optimisation --- monthly or quarterly, on 3--12 month trailing windows ---
has nothing real to track. The process parameters are homogeneous across the sample at
every horizon tested, including across the one genuine regime break available. A refit
scheme would be fitting estimation noise with 6--25 trades of signal per window, and
the 40-fold P\&L evidence already shows how that ends. This closes the question at the
level of the price process, which is stronger than closing it at the level of one more
walk-forward: there is no drift for \emph{any} P\&L refit scheme to capture, whatever
its cadence and window.

What the market's monthly variation actually calls for is already deployed (range
scaling, rolling AR(1), rolling residual sigma). If anything further is pursued, the
supported direction is sizing, not signal: the OU residual sigma trajectory is the
one series that moves with the regime, and exposure or buffer conditioning on it would
be a structural change --- the category that has survived out-of-sample before.
Reproduction: \texttt{premarket\_study/rolling\_mle.py} (study),
\texttt{test\_rolling\_mle.py} (instrument validation),
\texttt{mle\_homogeneity.py} (LR tests), \texttt{analyze\_mle\_drift.py} (tables and
figure); data extracted from the live workbook's Query sheet, not committed.

\end{document}
